% GENERATED FILE. Edit canonical lessons, then run npm run generate:books.
% canonical-source-hash: fnv1a64:38588f262563e737
% canonical-lessons: RU-C27-words, RU-C27-notices, RU-C27-pervoe-chtenie

\chapter{Reading — Six Words, Six Lines, and a First Passage}
\label{ch:ru-pervoe-chtenie}
\hlchaptermodality{27}

\begin{chapteropening}
\textbf{By the end of this chapter:} \emph{I can read Russian words off a board, read the lines I would say to a stranger, and hold a whole passage without translating it word by word.}

The introduction the book has been saying since chapter 2, read in order for the first time, with no new word in it.
\end{chapteropening}

\section[voda, khleb, chai, kofe, sup, syr]{(reading) — six words off a board}

\label{lesson:RU-C27-words}



\begin{quote}
You have written every letter in every one of these. Not one word is new.

What is new is that nobody says them before you read them.
\end{quote}

\begin{tcolorbox}[breakable,skin=enhanced,colback=orange!6,colframe=orange!45!black,boxrule=0.5pt,arc=1mm,left=6pt,right=6pt,top=4pt,bottom=4pt,fonttitle=\bfseries,title={Read this}]
\begin{quote}
\ru{вода} \ru{хлеб} \ru{чай} \ru{кофе} \ru{суп} \ru{сыр}
\end{quote}

Read down the list once. Do not spell each one out — look at the whole word and let it arrive.

Again, and notice which came at once and which you had to spell. Both are fine. The difference is what practice removes.
\end{tcolorbox}

\subsection*{What to know first}
These six would be on a board over a counter, and they are the six you would want first.

Two of them are traps of exactly the kind this book warned you about. \textbf{\ru{суп}} and \textbf{\ru{сыр}} look like \emph{cyn} and \emph{cbip} to an eye trained on English, and they are \emph{sup} and \emph{syr}.

That is the whole reason to read a word cold rather than after hearing it. Heard first, you never test whether the letters actually reached you.

\subsection*{Your turn}
\begin{itemize}
\raggedright
  \item \emph{Say it:} the six words down the list, once, without stopping
\end{itemize}

\emph{Twice through:}

\begin{tcolorbox}[breakable,colback=teal!4,colframe=teal!35!black,title={Before you move on}]
How many of the six were new? (\textbf{None.})

Which two look most like English and are not? (\textbf{\ru{суп}} and \textbf{\ru{сыр}}.)
\end{tcolorbox}
\section[privet, spasibo, pozhaluysta …]{(reading) — three to open, three to rescue}

\label{lesson:RU-C27-notices}



\begin{quote}
The last lesson gave you six things on a board. These are six things you would say.

Nothing here is new. What is new is meeting them in writing.
\end{quote}

\begin{tcolorbox}[breakable,skin=enhanced,colback=orange!6,colframe=orange!45!black,boxrule=0.5pt,arc=1mm,left=6pt,right=6pt,top=4pt,bottom=4pt,fonttitle=\bfseries,title={Read this}]
\begin{quote}
\ru{привет} \ru{спасибо} \ru{пожалуйста} \ru{извините} \ru{медленно} \ru{помогите}
\end{quote}

Read down once, without stopping.

Again — and notice the list turns over halfway. The first three open a conversation. The last three rescue one.
\end{tcolorbox}

\subsection*{What to know first}
\textbf{\ru{пожалуйста}} is the longest word on the list and the one you will read most often. It is worth recognising by shape rather than by spelling it out, because it does two jobs: \emph{please}, and \emph{you're welcome}.

\textbf{\ru{помогите}} and \textbf{\ru{извините}} end the same way, and both are requests you make to somebody you have not met. That shared ending is the one piece of grammar visible on this page.

Six lines, and half of them exist for the moment something goes wrong.

\subsection*{Your turn}
\begin{itemize}
\raggedright
  \item \emph{Say it:} the six lines, once, in order, without stopping
\end{itemize}

\emph{Twice through:}

\begin{tcolorbox}[breakable,colback=teal!4,colframe=teal!35!black,title={Before you move on}]
How many of the six were new? (\textbf{None.})

Where does the list turn over? (\textbf{After \ru{пожалуйста}} — three to open, three to rescue.)
\end{tcolorbox}
\section[privet! menya zovut Anna.]{(reading) — the first passage}

\label{lesson:RU-C27-pervoe-chtenie}



\begin{quote}
Six words, then six lines. Now the whole thing at once.

Not one word in it is new.
\end{quote}

\begin{tcolorbox}[breakable,skin=enhanced,colback=orange!6,colframe=orange!45!black,boxrule=0.5pt,arc=1mm,left=6pt,right=6pt,top=4pt,bottom=4pt,fonttitle=\bfseries,title={Read this}]
\begin{quote}
\ru{Привет}! \ru{Меня} \ru{зовут} \ru{Анна}. \ru{Сестра} \ru{и} \ru{мама}. \ru{Папа} \ru{тоже}. \ru{Я} \ru{читаю}. \ru{Я} \ru{хочу} \ru{читать}. \ru{Вода}, \ru{хлеб} \ru{и} \ru{чай}. \ru{Кофе} \ru{или} \ru{чай}? \ru{Я} \ru{не} \ru{знаю}. \ru{Спасибо}. \ru{До} \ru{свидания}.
\end{quote}

Read it through once without stopping. Do not translate as you go.

Again — and this time notice that two lines differ by one word, and that the one word changes when the reading happens.
\end{tcolorbox}

\subsection*{What to know first}
\textbf{\ru{Я} \ru{читаю}} is now. \textbf{\ru{Я} \ru{хочу} \ru{читать}} is not yet. Russian marks that difference by changing the ending of the verb in the first and leaving it whole in the second, which is the opposite of what an English speaker expects.

Look at \textbf{\ru{и}} and \textbf{\ru{или}} as well. One letter apart on the page, and the difference between a list and a choice. Read quickly, that is the easiest mistake on this page to make.

This is the first whole passage the book has asked you to hold at once, and it works because every line was paid for.

\subsection*{Your turn}
\begin{itemize}
\raggedright
  \item \emph{Say it:} the passage aloud, once, without stopping
\end{itemize}

\emph{Twice through:}

\begin{tcolorbox}[breakable,colback=teal!4,colframe=teal!35!black,title={Before you move on}]
How many words in that passage were new? (\textbf{None.})

Which two words are one letter apart and mean different things? (\textbf{\ru{и}} and \textbf{\ru{или}}.)
\end{tcolorbox}
